% COGENT: Cognitive Dimension-Aware Memory for Enhancing Code Agents
% 中文学术论文 — 参考 AAAI 2027 格式排版
% 注意：AAAI 禁止正文中使用非罗马字符，此中文版仅供内部参考，
% 正式投稿请使用 paper_en.tex

\documentclass[UTF8,twocolumn]{ctexart}

\usepackage{graphicx}
\usepackage{amsmath}
\usepackage{amssymb}
\usepackage{booktabs}
\usepackage{algorithm}
\usepackage{algorithmic}
\usepackage{newfloat}
\usepackage{listings}
\DeclareCaptionStyle{ruled}{labelfont=normalfont,labelsep=colon,strut=off}
\lstset{%
	basicstyle={\footnotesize\ttfamily},
	numbers=left,numberstyle=\footnotesize,xleftmargin=2em,
	aboveskip=0pt,belowskip=0pt,
	showstringspaces=false,tabsize=2,breaklines=true}
\floatstyle{ruled}
\newfloat{listing}{tb}{lst}{}
\floatname{listing}{Listing}
\usepackage{multirow}
\usepackage{array}
\usepackage{xcolor}
\usepackage{enumitem}
\usepackage[a4paper, margin=1in]{geometry}

\frenchspacing

\title{\textbf{COGENT: 面向代码智能体的认知维度感知记忆框架}}

\author{
    匿名作者\textsuperscript{\rm 1}
}
\date{}

\begin{document}
\maketitle

% ======================================================================
\begin{abstract}
基于大语言模型（LLM）的代码智能体通过工具调用轨迹求解编程任务，但每个问题均被独立处理——先前积累的操作性知识在回合切换时全部丢弃。近年来，智能体记忆研究 \cite{expel, reflexion, memgpt} 和层级化记忆架构 \cite{g-memory, evolve-mem} 已证明保留历史经验的价值，但现有系统或按时间顺序、或按嵌入距离、或按抽象层级单一维度来组织记忆，均忽略了每条经验所编码的\emph{推理类型}差异。本文提出 \textsc{COGENT}，一个按认知维度（因果、对比、策略、环境）与抽象层级（轨迹、工作流、摘要、洞察）双轴交叉组织代码智能体经验的记忆框架，构成 $4 \times 4$ 的结构化记忆网格。对每个新任务，\textsc{COGENT} 首先提取其认知画像以诊断所需的推理辅助类型，随后通过双通道相似度检索（任务级和认知级嵌入并行匹配），由维度原生 LLM 重排序器按各记忆自身认知类别的评分标准进行独立打分，最终组装互补的记忆组合。全部记忆通过系统提示注入，无需任何模型训练。在 SWE-bench Verified 的 Django 子集 100 个重叠任务上，以 DeepSeek-V4-Flash 为基座模型，\textsc{COGENT} 将通过率从 56.0\%（无记忆基线）提升至 76.0\%，净增 20 个百分点（22 个任务新解决，2 个退化）。我们进一步设计了包含 20 个实验条件的系统消融方案——涵盖核心流水线消融、认知维度逐一移除、抽象层级逐一移除、鲁棒性检验和多模型验证——以量化各组件的边际贡献，所有实验均已完成实现并处于执行阶段。
\end{abstract}

\begin{links}
    % \link{代码}{https://github.com/example/cogent}
    % \link{数据集}{https://example.com/datasets}
    % \link{扩展版本}{https://arxiv.org/abs/XXXX.XXXXX}
\end{links}

% ======================================================================
\section{引言}

以 mini-swe-agent、SWE-Agent 和 OpenHands 为代表的 LLM 代码智能体已能自主完成仓库导航、源文件检视和真实缺陷修复 \cite{swe-agent, mini-swe-agent, openhands}。其标准流程已然成熟：接收 issue 描述，在信息收集和代码编辑间交替迭代，补丁通过测试套件即为成功。

但一个结构性的低效问题始终存在。在一次任务中积累的诊断推理链条、被探索后又放弃的死胡同、连接症状与根因的因果链路——这些操作性知识在回合之间完全蒸发。一个花费三十步定位出 Django 迁移顺序依赖违规的智能体，面对下一个因果结构几乎完全相同的迁移问题时，仍需从头开始：重列目录、重读配置、重绘依赖图。

这一观察催生了大量智能体记忆研究。Reflexion \cite{reflexion} 维护来自失败尝试的口头自我反思。ExpeL \cite{expel} 通过比较成功与失败轨迹来系统性地提取跨任务洞察，以 ADD/EDIT/UPVOTE/DOWNVOTE 算子管理不断演化的洞察池。层级化记忆架构——G-Memory 的三层图结构 \cite{g-memory}、EVOLVE-MEM 的自适应层级 \cite{evolve-mem}、H-MEM 的四级位置索引 \cite{h-mem}——将经验组织为多个抽象粒度以实现高效的自顶向下检索。这些进展共同确立了"保留并复用智能体经验可改善下游任务表现"这一基本结论。

然而，现有记忆设计共享一个组织原则：记忆按照它们描述了\emph{什么}（哪个任务、哪个仓库、哪个抽象层级）来结构化，而非它们编码了\emph{哪种推理}。一条关于诊断过期缓存失效的记忆与一条 Django 会话超时的查询在词元层面几乎毫无重叠，但二者需要完全相同的认知操作——追踪写入路径、检查失效是否在正确信号上触发、验证键命名空间。反过来说，两条以高度相似语言描述同一 Django 子系统的记忆，可能编码着截然不同的推理类型（一个涉及模板转义，另一个涉及数据库事务边界）。表层文本邻近度是认知适用性的充满噪声的代理变量。

\textsc{COGENT} 通过引入\emph{认知维度轴}作为一等组织原则来偏离上述范式。每条记忆不仅标注其抽象层级，还标注其认知维度——因果（错误为何发生？修复为何有效？）、对比（哪些方案被尝试并放弃了？方案在什么边界条件下成立？）、策略（做出了什么工作流设计选择？理由是什么？）和环境（遇到了哪些仓库特定约定、工具链陷阱和配置细节？）。将这四个认知维度与四个抽象层级交叉，形成 $4 \times 4$ 的记忆网格，其中 16 个位置各自存放一类独特的操作性知识。

检索端同样体现这一结构。面对新任务时，\textsc{COGENT} 首先提取\emph{认知画像}，诊断该问题需要哪些推理维度的辅助。该画像驱动两个并行嵌入通道：任务级通道匹配具体语义（文件名、错误信息、函数签名），认知级通道匹配抽象推理模式（因果链拓扑、反模式签名）。随后，候选记忆被\emph{维度原生重排序}：每条记忆按其自身认知维度对应的评分标准被独立评估（因果记忆按因果链同构度评分，环境记忆按仓库区域重叠度评分），被画像判定为低需求的维度获得惩罚系数。最后，贪心选择器组装最终记忆集合，惩罚近似重复对，奖励互补维度组合。

\textsc{COGENT} 完全通过提示注入运行——所有选中的记忆在执行前被拼接入智能体的系统消息。这种黑盒设计不需要任何参数更新、架构修改或与特定 LLM 的耦合，可立即部署于任何通过 API 访问的代码智能体。

本文贡献如下：

\begin{enumerate}[label=(C\arabic*), leftmargin=*]
    \item 提出面向代码智能体记忆的认知维度分类体系——因果、对比、策略、环境——按推理类型而非表面内容来归类经验，并将其与抽象层级交叉构成 $4 \times 4$ 记忆网格，形成有别于纯层级化或纯扁平化存储的双轴组织结构。
    \item 设计了一套检索协议，包含认知画像提取（检索前诊断查询的推理需求）、双通道相似度匹配（任务级和认知级嵌入加权融合）以及维度原生重排序（按候选记忆自身认知类别的评分标准进行独立打分）。
    \item 建立了增量序列记忆累积范式——记忆池逐轮扩展，每个已解决任务贡献其提取的网格至共享存储，从而能够在贴近真实持续部署的条件下测量单任务记忆的边际贡献。
    \item 设计了一套包含 20 个实验条件的系统消融方案——覆盖核心流水线消融、认知维度逐一移除、抽象层级逐一移除、检索鲁棒性检验和多模型验证——所有实验均已完成实现并处于执行阶段。
    \item 在 SWE-bench Verified 的 100 个 Django 子集任务上进行了实证验证：\textsc{COGENT} 达到 76.0\% 通过率，对比 56.0\% 的无记忆基线，净增 20 个百分点（22 个任务改善，2 个退化）。
\end{enumerate}

% ======================================================================
\section{相关工作}

\subsection{从智能体轨迹中提取经验}

从智能体执行痕迹中蒸馏可复用知识的思路已被广泛探索。Reflexion \cite{reflexion} 开创了口头的自我反思：失败后，智能体分析自身轨迹并生成自然语言诊断，在下一次尝试时前置到提示词中。ExpeL \cite{expel} 将此推广为跨任务洞察提取，利用 LLM 比较同一任务的成功与失败轨迹，以及识别不同任务间共享的成功模式，以显式的 ADD/EDIT/UPVOTE/DOWNVOTE 算子维护不断演化的洞察池。CTIM-Rover \cite{ctim-rover} 随后尝试将 ExpeL 的经验学习应用于软件工程智能体，但发现仓库级别的跨任务记忆引入了噪声，将智能体偏向无关的词元——这恰好凸显了\emph{选择性检索}的挑战，而我们的维度感知方法正是针对这一问题设计的。DoT-Bank \cite{dot-bank} 和 Agent Workflow Memory \cite{awm} 进一步推进了轨迹到记忆的蒸馏技术。

\textsc{COGENT} 建立在上述工作基础之上——我们同样从已完成的轨迹中提取结构化记忆——但在\emph{组织模式}上有根本性不同。ExpeL 将洞察存储为由投票机制管理的无结构自然语言池，Reflexion 维护的是每个任务的口头反思，而 \textsc{COGENT} 施加了一个二维认知结构（维度 $\times$ 层级），使得检索能够匹配查询的认知需求，而非仅依赖通用相似度。

\subsection{层级化记忆架构}

多层级记忆组织近期受到大量关注。G-Memory \cite{g-memory} 将多智能体交互轨迹结构化为三层图（洞察图 $\rightarrow$ 查询图 $\rightarrow$ 交互图），支持双向自顶向下/自底向上遍历。EVOLVE-MEM \cite{evolve-mem} 实现了包含动态记忆网络、层级化管理器和自我改进引擎的自适应层级结构。H-MEM \cite{h-mem} 提出了四级位置索引以实现高效路由。Agent-Memo 使用受文件系统启发的分层加载方案。这些系统展示了在多个粒度上组织记忆的价值，我们的抽象层级轴（轨迹 $\rightarrow$ 工作流 $\rightarrow$ 摘要 $\rightarrow$ 洞察）分享了这一动机。\textsc{COGENT} 的差异在于添加了第二个正交轴：认知维度。我们没有仅按抽象粒度组织，而是将抽象层级与推理类型相交叉，创建一个每个坐标代表一类独特有用知识的网格。

\subsection{记忆冗余控制与互补性检索}

多项近期工作致力于从候选池中选择互补、非冗余的信息。Context-Picker \cite{context-picker} 使用两阶段 GRPO 强化学习，第二阶段收紧冗余边界以剪枝干扰项。MEM1 \cite{mem1} 通过端到端 RL 学习协同记忆巩固与推理，策略性地丢弃冗余信息。MemFly \cite{memfly} 将记忆选择形式化为在线信息瓶颈问题，在最小化压缩熵的同时最大化相关性。我们的协同感知贪心选择——带显式成对冗余惩罚（相同层级、相同维度、高余弦相似度）和互补奖励（相同源任务、不同认知维度）——在同一问题空间中运作。相对于这些学习方法，我们基于规则的成对调整更简单、无需训练，但共享"记忆选择应优化联合覆盖而非独立逐项分数"这一核心洞察。

\subsection{认知科学基础}

\textsc{COGENT} 认知维度分类体系的设计受人类记忆研究的经典发现启发。Tulving 的情节-语义记忆二分法 \cite{tulving-episodic} 区分了对具体事件的记忆与对一般知识的记忆，催生了我们的抽象层级轴。Schank 的动态记忆理论 \cite{schank-dynamic} 主张表面不相似情境之间的"触发回忆"由共享的主题组织结构介导，直接支持了我们的核心主张——认知维度对齐可以浮现词汇搜索会遗漏的相关记忆。编码特异性原则 \cite{tulving-encoding} 为认知画像提取步骤提供了理论依据。近期的 LLM 智能体记忆系统包括 DCPM \cite{dcpm} 和 CDMem \cite{cdmem} 也借鉴了认知科学，但其组织原则（双过程理论、上下文依赖索引）与我们的维度分类体系不同。

% ======================================================================
\section{COGENT 框架}

我们从三个方面描述 \textsc{COGENT}：定义存储组织的认知记忆网格（第~\ref{sec:lattice-cn} 节）、为每个查询选择记忆的检索流水线（第~\ref{sec:retrieval-cn} 节）以及序列池协议（第~\ref{sec:pool-cn} 节）。从已完成轨迹构建记忆遵循已有的提取实践 \cite{expel, reflexion}：原始轨迹首被压缩为结构化 JSON（一次 LLM 调用），随后四次 LLM 调用——每个抽象层级一次——各自同时提取全部四个认知维度的条目（总计 16 个单元格、5 次 LLM 调用）。我们省略对此标准过程的详细描述，聚焦于\textsc{COGENT}的独特之处：认知网格结构及其上的检索机制。

% ----------------------------------------------------------------------
\subsection{认知记忆网格}
\label{sec:lattice-cn}

\textsc{COGENT} 的核心组织结构是由两条正交轴线定义的 $4 \times 4$ 网格。每条记忆条目 $m$ 占据由其抽象层级 $l$ 和认知维度 $d$ 确定的坐标 $(l, d)$。

\subsubsection{抽象层级轴（纵向）}

纵轴涵盖四个由具体到抽象的层级。类似的层级组织见于 G-Memory \cite{g-memory}、EVOLVE-MEM \cite{evolve-mem} 和 H-MEM \cite{h-mem}；我们的具体层级定义针对代码调试轨迹的结构进行了定制：

\begin{itemize}
    \item \textbf{轨迹层}（$L_1$）：命令级操作记录，保存执行的确切动作序列、工具输出、检视和修改的文件以及探索的分支结构。
    \item \textbf{工作流层}（$L_2$）：步骤级过程模式，对单条命令进行抽象，捕获子目标依赖、流程控制决策和调试阶段的时序组织。
    \item \textbf{摘要层}（$L_3$）：任务级诊断综合，识别根因机制、附带理由的解决策略以及从观察症状到底层缺陷的映射。
    \item \textbf{洞察层}（$L_4$）：跨任务方法论原则——启发式规则、设计模式、调试策略——以不引用特定仓库或实例的方式表述。
\end{itemize}

\subsubsection{认知维度轴（横向）}

横轴是 \textsc{COGENT} 的核心概念贡献。它按照记忆所编码的\emph{推理类型}进行分类，独立于其抽象层级：

\begin{itemize}
    \item \textbf{因果维度}：关于原因和结果的诊断推理。每条因果记忆包含 \texttt{error\_category}、\texttt{cause\_category}、\texttt{intervention\_type} 和结构化的 \texttt{causal\_chain}。这些记忆回答\emph{为什么}的问题——为什么错误会出现、为什么所选修复能恢复正确行为、为什么替代方案被拒绝。
    \item \textbf{对比维度}：区分有效与无效方法的比较知识。包含\emph{反模式}（被尝试但失败的操作及失败原因）、\emph{成功模式}（成功的操作及区分性特征）和\emph{边界条件}（成功方法的适用范围和失效条件）。这些记忆回答\emph{什么区分了}成功与失败。
    \item \textbf{策略维度}：关于规划与过程的元认知知识——在备选方案间选择的决策准则、工作流设计的理由、风险评估启发式规则以及组织调试过程的指导方针。这些记忆回答\emph{如何着手}的问题。
    \item \textbf{环境维度}：仓库特定和工具链特定的操作知识——目录布局惯例、依赖图结构、测试框架调用模式、常见配置陷阱（如迁移排序约束、fixture 加载语义）和工具调用惯用法。这些记忆回答在特定代码库中\emph{在哪里和如何操作}的问题。
\end{itemize}

这四个维度的选择基于对 SWE-bench 上代码智能体失败模式的分析。对失败轨迹的检视表明，智能体错误可聚类为四类：(i) 误判连接症状与根因的因果机制（因果），(ii) 坚持一个表面合理但最终错误的修复策略（对比），(iii) 遵循低效或误导的探索顺序（策略），(iv) 未能发现仓库特定约定或工具链行为（环境）。四个维度直接针对这些观察到的失败类别。

\subsubsection{记忆条目结构}

网格中每个单元存储三类信息。

\textbf{认知摘要。} 对记忆所捕获推理内容的一份简短、去情境化描述——一段因果链、一个反模式、一条工作流策略、一项仓库惯例——在撰写时刻意不提及任何具体文件、模块或函数名。系统为记忆生成嵌入时，供给认知级嵌入通道的正是这份摘要（而非完整任务文本）。LLM 重排序候选时，阅读的也是这些摘要，从而在不受表层细节干扰的情况下评判认知匹配度。

\textbf{具体锚点。} 锚点与摘要相反：刻意具体、刻意局部。它们记录原始轨迹中出现的实际文件路径、函数名、错误信息和配置键。推理时，智能体可将检索记忆的锚点与当前正在检视的文件和函数做对比。锚点指向同一子系统的记忆很可能具体适用；锚点毫无交集的记忆可能在认知上相关但实践上错位。

\textbf{复合嵌入文本。} 我们将任务描述、认知摘要和具体锚点拼接为一段文本，用 SentenceTransformer（all-MiniLM-L6-v2，384 维向量）\cite{sentence-transformers} 编码。从这一段文本中提取两组嵌入。第一组——称为任务键嵌入——保留仓库特定的词汇：错误类型、API 名称、文件路径。第二组——认知嵌入——以摘要字段为主导，捕捉基本独立于表层词汇的推理结构。下文阶段~2 所述的双通道检索即运行于这对向量之上。

% ----------------------------------------------------------------------
\subsection{检索流水线}
\label{sec:retrieval-cn}

给定查询任务 $q$，检索流水线选择 $K=3$ 条互补记忆用于提示注入。流水线最多消耗 3 次 LLM 调用。

\subsubsection{阶段 1：认知画像提取（1 次 LLM 调用）}

第一阶段刻画查询需要\emph{哪些类型}的推理辅助。LLM 分析任务描述并输出认知画像：

\begin{equation}
    P(q) = \{\text{need\_causal}, \text{need\_contrastive}, \text{need\_strategic}, \text{need\_environment}\} \in [0,1]^4
\end{equation}

同时提取因果签名 $\Sigma(q) = \{\text{error\_category}, \text{cause\_category}, \text{intervention\_type}, \text{causal\_chain}\}$，假设性地描述错误的因果结构。画像有两个用途：为认知级嵌入通道提供文本（阶段 2），以及通过标记低需求维度来控制维度原生重排序的惩罚机制（阶段 3）。

\subsubsection{阶段 2：双通道语义检索（0 次 LLM 调用）}

两个独立的嵌入相似度通道计算每个记忆 $m_i$ 的相关性估计：

\begin{align}
    \text{task\_sim}_i &= \cos\big(\text{emb}(q_{\text{task}}),\; \text{emb}_{\text{key}}(m_i)\big) \\
    \text{cog\_sim}_i &= \cos\big(\text{emb}(q_{\text{cog}}),\; \text{emb}_{\text{cog}}(m_i)\big)
\end{align}

通道一（$\text{task\_sim}$）将原始任务描述与复合嵌入文本匹配，捕获仓库词汇、错误签名和函数名的具体重叠。通道二（$\text{cog\_sim}$）将认知画像文本与纯认知摘要匹配，捕获独立于表面词汇的推理结构相似性。两通道融合：

\begin{equation}
    \text{semantic}_i = \alpha_{\text{dual}} \cdot \text{task\_sim}_i + (1 - \alpha_{\text{dual}}) \cdot \text{cog\_sim}_i
\end{equation}

其中 $\alpha_{\text{dual}} = 0.70$，优先具体相似度，同时让认知通道恢复结构同构但词汇疏远的候选。按 $\text{semantic}_i$ 排名前 $N=20$ 的条目进入重排序。

\subsubsection{阶段 3：维度原生重排序（1 次批量 LLM 调用）}

此阶段是 \textsc{COGENT} 的核心检索创新。我们不是让 LLM 为每个候选生成一个通用相关性分数，而是按照\emph{候选记忆自身认知维度对应的标准}来评估每条记忆：

\begin{itemize}
    \item \textbf{因果}条目按其存储的 $\texttt{causal\_chain}$ 与查询假设的 $\Sigma(q)$ 之间的结构对齐程度评分——错误$\to$原因$\to$修复的骨架是否匹配？
    \item \textbf{对比}条目按其反模式或成功模式是否与查询任务可能展现的失败模式相关来评分。
    \item \textbf{策略}条目按其方法论指导对查询的仓库上下文和问题结构的可迁移性来评分。
    \item \textbf{环境}条目按其描述的仓库区域与查询涉及的文件或子系统之间的重叠度来评分。
\end{itemize}

对于 $P(q)$ 标记为低需求的维度（$\text{need\_d} < 0.3$），维度分数乘以 0.65 的惩罚系数。该惩罚机制可操作化了以下直觉：即使一条因果记忆评分很高，若查询主要需要环境特定的操作知识也于事无补。

$N=20$ 条候选的认知摘要被打包为单次批量 LLM 调用以最小化 token 成本。融合最终分数为：

\begin{equation}
    \text{combined}_i = \alpha_s \cdot \text{semantic}_i + \alpha_c \cdot \text{cognitive}_i
\end{equation}

其中 $\alpha_s = 0.35$，$\alpha_c = 0.65$。认知项权重更高体现了我们的设计原则：语义通道提供宽泛候选池后，认知结构匹配度应主导最终排序。

\subsubsection{阶段 4：成对交互感知组装（0 次 LLM 调用）}

最终阶段在选择 $K=3$ 条记忆时考虑集合内部的成对关系。设 $\mathcal{S}$ 为已选记忆集合。第 1 轮选择 $m^* = \arg\max_i \text{combined}_i$。后续轮次中，每个剩余候选 $C$ 与每个已选 $S \in \mathcal{S}$ 的成对交互计算如下：

\begin{itemize}
    \item \textbf{冗余惩罚}：当 $C$ 与 $S$ 属同一抽象层级、同一认知维度、且余弦相似度超过 0.92 时，施加 $\delta_{\text{red}} = -0.15$。
    \item \textbf{互补奖励}：当 $C$ 与 $S$ 来自同一源任务但属不同认知维度时，施加 $\delta_{\text{comp}} = +0.10$。因果+环境配对额外给予增强奖励 $+0.15$，二者共同提供诊断推理和具体的操作知识。
\end{itemize}

调整分数为 $\text{adjusted}(C) = \text{combined}(C) + \sum_{S \in \mathcal{S}} (\delta_{\text{comp}}(C, S) - \delta_{\text{red}}(C, S))$。最高调整分数的候选进入 $\mathcal{S}$。

我们采用简单的基于规则的公式化而非学习的冗余控制 \cite{context-picker, memfly}，原因有二：无需训练数据，且成对交互极为稀疏（每次选择至多 $K \times (K-1)/2 = 3$ 次比较），基于规则的方法在计算上微不足道且保持可解释性。

\subsubsection{自适应分数阈值}

组装完成后，对选中记忆施加一道批量自适应的质量过滤。阈值 $\tau$ 基于全部 $N=20$ 条候选的分数分布（而非仅已选的 $K=3$ 条）计算：$\tau = \max(\tau_{\text{floor}},\; \text{median}(\{\text{combined}_i\}) - \lambda \cdot \sigma(\{\text{combined}_i\}))$，其中 $\tau_{\text{floor}} = 0.45$，$\lambda = 0.5$。使用 20 条的完整分布是为了获得可靠的统计量——3 个样本不足以进行有意义的均值-标准差计算。该阈值随后应用于已选记忆：组合分数低于 $\tau$ 者被丢弃，始终保留至少一条记忆。当候选批次整体质量较高时阈值自动收紧（鼓励选择性），当分数普遍偏低时阈值回落至地板值（避免在低质量批次中丢弃全部记忆）。

% ----------------------------------------------------------------------
\subsection{序列记忆池协议}
\label{sec:pool-cn}

\textsc{COGENT} 支持两种运作模式。\emph{批量}模式下，一组固定的源任务填充一个静态记忆池 $\mathcal{M}_{\text{batch}}$ 供所有目标查询使用，适用于受控离线评估。\emph{序列}模式——本文的主要评估协议——任务按固定顺序逐一处理，池在每个回合后增量扩展：

\begin{equation}
    \text{第 } t \text{ 轮}: \quad \text{Task}_t \xrightarrow{\text{从 } \{\mathcal{M}_0, \ldots, \mathcal{M}_{t-1}\} \text{ 检索}} \text{执行} \xrightarrow{\text{提取}} \mathcal{M}_t \rightarrow \text{Pool}_{1:t}
\end{equation}

每个新解决的任务将其提取的网格贡献至共享存储，立即可供后续轮次使用。该协议模拟持续部署场景：智能体随时间累积经验，每个回合受益于全部历史。至关重要的是，序列模式可直接测量每个任务的\emph{边际贡献}——通过追踪池从 1 增长至 99 的过程中累计通过率的演化，观察增量记忆价值如何变化、饱和点何时出现。

% ======================================================================
\section{实验设计}

以下所有实验均已完成实现并处于执行阶段。我们报告完整的设计方案；完整系统（E1）和无记忆基线（E6）的结果来自已完成的运行，E2--E5、E7--E17 和 E25--E26 的结果在运行尚未完成处标记为待定。\footnote{所有实验的代码和配置文件见项目仓库 \texttt{harbor/configs/experiments/}。}

\subsection{实验配置}

\textbf{数据语料。} 从 SWE-bench Verified 的 Django 子集 \cite{swe-bench} 中以种子 42 随机采样 100 个重叠任务。每个任务包含自然语言 issue 描述和用于自动验证的预留测试套件。

\textbf{智能体与模型。} mini-swe-agent 为基础智能体，DeepSeek-V4-Flash 为后端 LLM（\texttt{deepseek-chat}，温度 0）。嵌入使用 SentenceTransformer 的 all-MiniLM-L6-v2 模型（384 维）。所有实验共享相同模型配置，智能体每任务最多 150 步。无记忆基线的系统提示仅含标准智能体指令和任务描述。

\textbf{评估。} 主要指标：pass@1（每任务单次尝试）。次要指标：基线失败$\rightarrow$COGENT 成功的任务数（改善）和基线成功$\rightarrow$COGENT 失败的任务数（退化）。序列实验额外追踪累计通过率随池大小的变化。

\textbf{完整系统配置（E1）。} $\alpha_{\text{dual}} = 0.70$，$N = 20$，$\alpha_s = 0.35$，$\alpha_c = 0.65$，$K = 3$，$\tau_{\text{floor}} = 0.45$，$\lambda = 0.5$，冗余余弦阈值 $0.92$。所有实验采用序列执行协议，任务排序一致（随机种子 42），确保跨实验可比性。

\subsection{核心流水线消融阶梯（E1--E6）}

表~\ref{tab:ablation-config-cn} 定义了一个六实验消融阶梯，从完整系统逐步剥离组件，回答关于各组件边际价值的嵌套问题。

\begin{table}[t]
\centering
\caption{E1--E6 的配置矩阵。共享条件：序列模式、100 任务（种子 42）、mini-swe-agent、DeepSeek-V4-Flash、步数限制 150。}
\label{tab:ablation-config-cn}
\begin{tabular}{lccccc}
\toprule
\textbf{ID} & \textbf{记忆结构} & \textbf{重排序} & \textbf{协同} & \textbf{通道} & \textbf{阈值} \\
\midrule
E1 (完整)   & $4{\times}4$ 认知网格 & \checkmark & \checkmark & 双通道 & 自适应 \\
E2          & $4{\times}4$ 认知网格 & ---        & \checkmark & 双通道 & 自适应 \\
E3          & $4{\times}4$ 认知网格 & \checkmark & ---        & 双通道 & 自适应 \\
E4          & $4{\times}4$ 认知网格 & ---        & ---        & 单通道 & 固定 \\
E5          & MTL 三层（无维度）   & ---        & ---        & 单通道 & 固定 \\
E6          & 无（基线）          & ---        & ---        & ---    & --- \\
\bottomrule
\end{tabular}
\end{table}

阶梯对应五个嵌套比较：

\begin{itemize}
    \item \textbf{E5 vs.\ E6}：\emph{任何}形式的记忆——即使没有认知维度结构——是否优于完全无记忆？量化"记忆本身"的价值。
    \item \textbf{E4 vs.\ E5}：将扁平三层方案替换为 $4 \times 4$ 认知网格（同时保持检索机制不变——纯嵌入、无 LLM 增强）的增量收益是多少？隔离双轴认知组织的表征贡献。
    \item \textbf{E2 vs.\ E1}：维度原生 LLM 重排序（阶段 3）在双通道嵌入之上提供了多少额外增益？隔离维度原生评分标准的价值。
    \item \textbf{E3 vs.\ E1}：成对交互感知组装（阶段 4）的独立贡献是多少？隔离对已选集合内显式冗余/互补建模的价值。
    \item \textbf{E4 vs.\ E1}：从纯嵌入到完整流水线的全跨度——基于 LLM 的认知组件（重排序 + 交互感知选择）共同带来的总增益是多少？
\end{itemize}

\subsection{认知维度逐一移除（E7--E10）}

为验证每个认知维度贡献非冗余信息，从完整 $4 \times 4$ 网格中逐一移除单个维度的全部记忆，保留完整 LLM 流水线：

\begin{itemize}
    \item \textbf{E7（--因果）}：检索池限定为对比 + 策略 + 环境。
    \item \textbf{E8（--对比）}：检索池限定为因果 + 策略 + 环境。
    \item \textbf{E9（--策略）}：检索池限定为因果 + 对比 + 环境。
    \item \textbf{E10（--环境）}：检索池限定为因果 + 对比 + 策略。
\end{itemize}

若移除某维度后通过率显著下降，即构成该维度编码了其余三个维度无法弥补的独特信息的证据。

\subsection{抽象层级逐一移除（E11--E13）}

平行分析从完整网格中逐一移除单个层级：

\begin{itemize}
    \item \textbf{E11（--轨迹层）}：检索池限定为工作流 + 摘要。
    \item \textbf{E12（--工作流层）}：检索池限定为轨迹 + 摘要。
    \item \textbf{E13（--摘要层）}：检索池限定为轨迹 + 工作流。
\end{itemize}

\subsection{检索鲁棒性检验（E14--E17）}

\begin{itemize}
    \item \textbf{E14（随机记忆）}：用均匀随机采样 $K=3$ 条记忆替代整个检索流水线，确立检索质量的下界。
    \item \textbf{E15（LLM 直接 Top-K）}：绕过认知维度分类体系——将全部 $N=20$ 条候选（不带维度标签）呈现给 LLM，要求直接选择最佳 $K=3$。检验显式维度结构是否提供 LLM 隐式判断之外的额外价值。
    \item \textbf{E16（固定阈值）}：禁用自适应阈值（$\lambda = 0$），使用恒定阈值 $\tau = 0.45$。检验分布感知的动态阈值是否提供可测量的收益。
    \item \textbf{E17（无洞察附加）}：抑制将抽象洞察层记忆链接到被选中具体记忆的机制。测量去情境化的方法论原则在具体操作知识之外的增量价值。
\end{itemize}

\subsection{多模型验证（E25--E26）}

使用 GPT-4o（E25）和 Claude Opus 4.5（E26）作为后端模型复现完整系统配置（E1）。

\subsection{主要结果（已完成）}

\begin{table}[t]
\centering
\caption{SWE-bench Verified Django 子集（100 任务）的主要结果。}
\label{tab:main-results-cn}
\begin{tabular}{lccc}
\toprule
\textbf{条件} & \textbf{通过率} & \textbf{改善} & \textbf{退化} \\
\midrule
无记忆（基线，E6） & 56.0\% & --- & --- \\
\textsc{COGENT}（完整，E1） & \textbf{76.0\%} & 22 & 2 \\
\bottomrule
\end{tabular}
\end{table}

表~\ref{tab:main-results-cn} 呈现主要对比。20 个百分点的净增益分解为 22 个新解决的任务和 2 个退化。改善案例涵盖多种重复出现的缺陷类型：数据库迁移排序违规（7 例）、模板渲染逻辑错误（5 例）、URL 分发与中间件配置错误（4 例）、表单验证边界情况（3 例）和序列化不一致（3 例）。收益集中在两类任务：(i) 因果结构在记忆池中有强先例的任务（如迁移依赖违规，其错误$\to$原因$\to$修复链条高度模式化），(ii) 环境知识密集型任务（智能体原本需花费大量步骤重新发现操作细节，如测试调用惯用法、fixture 加载行为）。

\subsection{消融结果（待填充）}

\begin{table}[t]
\centering
\caption{消融结果框架。标记为"待测"的单元格对应正在执行中的实验。}
\label{tab:ablation-results-cn}
\begin{tabular}{lcc}
\toprule
\textbf{条件} & \textbf{通过率} & \textbf{相对 E1 的 $\Delta$} \\
\midrule
E1: 完整 \textsc{COGENT} & 76.0\% & --- \\
E2: 无维度原生重排序 & 待测 & 待测 \\
E3: 无成对交互组装 & 待测 & 待测 \\
E4: 纯嵌入（无 LLM 阶段） & 待测 & 待测 \\
E5: MTL 三层（无认知维度） & 待测 & 待测 \\
E6: 无记忆 & 56.0\% & $-20.0$ pp \\
\midrule
\multicolumn{3}{c}{\textit{维度逐一移除}} \\
\midrule
E7: 去除因果 & 待测 & 待测 \\
E8: 去除对比 & 待测 & 待测 \\
E9: 去除策略 & 待测 & 待测 \\
E10: 去除环境 & 待测 & 待测 \\
\midrule
\multicolumn{3}{c}{\textit{层级逐一移除}} \\
\midrule
E11: 去除轨迹层 & 待测 & 待测 \\
E12: 去除工作流层 & 待测 & 待测 \\
E13: 去除摘要层 & 待测 & 待测 \\
\midrule
\multicolumn{3}{c}{\textit{鲁棒性检验}} \\
\midrule
E14: 随机记忆 & 待测 & 待测 \\
E15: LLM 直接 Top-K（无维度） & 待测 & 待测 \\
E16: 固定阈值 & 待测 & 待测 \\
E17: 无洞察附加 & 待测 & 待测 \\
\midrule
\multicolumn{3}{c}{\textit{多模型}} \\
\midrule
E25: GPT-4o（完整流水线） & 待测 & 待测 \\
E26: Claude Opus 4.5（完整流水线） & 待测 & 待测 \\
\bottomrule
\end{tabular}
\end{table}

\subsection{退化案例分析}

两个任务在引入记忆后出现退化。

\textbf{\texttt{django-13807}}。核心问题涉及 Django 的 \texttt{refresh\_from\_db()} 与用户自定义 \texttt{@property} 之间的交互，后者遮蔽了同名字段。检索到的因果记忆关于模型字段访问模式评分很高，但描述的是\emph{延迟属性加载}——与属性遮蔽在结构上邻近但因果上截然不同的机制。智能体在识别出真正的属性遮蔽问题前，花费了若干步骤探查延迟加载基础设施。这暴露了因果结构匹配的局限：当两个现象共享相似的因果表层签名（"访问模型属性返回了意料之外的数据"）但底层机制不同时，即使维度原生评分也可能产生假阳性。在因果维度内引入更细粒度的\emph{机制级}区分可能抑制此类错误。

\textbf{\texttt{django-14053}}。Django admin 模板标签注册中的命名空间冲突，实际只需在特定注册函数中进行三行局部修改。检索到的策略记忆强调"封装模板渲染逻辑"，一条对比反模式警告"避免全局模板状态修改"——它们将智能体引向了过度的重构方案。这暴露了策略层记忆的一种失效模式：抽象的方法论原则可能使智能体偏向过度工程化的解决方案，而真正的修复是小型且上下文特定的。

\subsection{计划中的分析}

待消融套件完成后，我们将进行以下分析：

\textbf{累计通过率随池大小的变化。} 将第 $1 \ldots t$ 轮累计通过率对池大小作图，揭示增量记忆的边际价值。预计呈现三个阶段：早期快速提升（池大小 1--15，每个新任务贡献大量新颖的认知覆盖）、增长趋缓（16--50，填补日益细分的认知缝隙）和近饱和（超过约 50，额外条目主要复现已有模式）。该轨迹将与人类记忆研究中的认知负荷饱和曲线 \cite{cognitive-load} 相呼应——检索干扰最终抵消额外存储信息的收益 \cite{ctim-rover}。

\textbf{维度选择频率。} 汇总所有任务中 top-$K$ 记忆的认知维度标签，刻画系统在经验上偏好哪些类型的辅助。初步检查表明因果记忆占比最高（诊断推理在调试中的中心性），其次为环境（操作知识）、策略（方法论）和对比（反模式识别——可能在智能体主动探索错误方向时选择性发挥价值）。

% ======================================================================
\section{结论}

本文提出 \textsc{COGENT}，一个围绕认知维度分类体系（因果、对比、策略、环境）与抽象层级交叉构建的代码智能体记忆框架。检索流水线从新查询中提取认知画像，执行双通道嵌入搜索，以维度原生标准重排序候选记忆，并通过对成对交互的考量组装互补记忆集合。框架完全通过提示注入运行，无需任何模型训练。

\textbf{局限性。} (1) 记忆构建每任务产生 5 次固定 LLM 调用开销，在超大规模部署时可将提取在结构相似的任务批次间摊销。(2) 认知维度分类体系为先验定义；基于数据的维度发现是否能产生更有效的划分仍为开放问题。(3) 当前评估限定于单一仓库生态（Django）和单一智能体框架（mini-swe-agent），跨仓库、跨智能体的泛化需专门研究。(4) 退化案例表明，当表面因果相似性掩盖底层机制差异时，高分认知匹配可能误导智能体，推动在因果维度内引入更细粒度的机制级区分。

\textbf{未来方向。} (1) 通过系统分析跨仓库和跨语言的智能体失败模式来实证优化维度清单；(2) 构建机制级因果区分器以抑制假阳性检索；(3) 探索可学习的、仓库自适应的维度权重；(4) 拓展评估至 AIME、LiveCodeBench 和 USACO；(5) 研究跨智能体知识迁移——由一个智能体架构提取的记忆能否使不同的智能体实现受益。

% ======================================================================
\section*{致谢}

% TODO

% ======================================================================
\begin{thebibliography}{}

\bibitem[Lewis et al., 2020]{rag-lewis}
Lewis, P.; Perez, E.; Piktus, A.; Petroni, F.; Karpukhin, V.; Goyal, N.; K\"{u}ttler, H.; Lewis, M.; Yih, W.; Rockt\"{a}schel, T.; Riedel, S.; and Kiela, D.
\newblock Retrieval-augmented generation for knowledge-intensive {NLP} tasks.
\newblock In \emph{NeurIPS}, 2020.

\bibitem[Yang et al., 2024]{swe-agent}
Yang, J.; Jimenez, C. E.; Wettig, A.; Lieret, K.; Yao, S.; Narasimhan, K.; and Press, O.
\newblock {SWE}-Agent: Agent-computer interfaces enable automated software engineering.
\newblock In \emph{NeurIPS}, 2024.

\bibitem[mini-swe-agent]{mini-swe-agent}
mini-swe-agent: A lightweight software engineering agent.
\newblock \emph{https://github.com/ethanmooney/mini-swe-agent}, 2024.

\bibitem[Wang et al., 2024]{openhands}
Wang, X.; et al.
\newblock {OpenHands}: An open platform for {AI} software developers as generalist agents.
\newblock \emph{arXiv preprint arXiv:2407.16741}, 2024.

\bibitem[Shinn et al., 2023]{reflexion}
Shinn, N.; Cassano, F.; Gopinath, A.; Narasimhan, K.; and Yao, S.
\newblock {Reflexion}: Language agents with verbal reinforcement learning.
\newblock In \emph{NeurIPS}, 2023.

\bibitem[Zhao et al., 2024]{expel}
Zhao, A.; Huang, D.; Xu, Q.; Lin, M.; Liu, Y.-J.; and Huang, G.
\newblock {ExpeL}: {LLM} agents are experiential learners.
\newblock In \emph{Proceedings of AAAI}, 2024.

\bibitem[Lindenbauer et al., 2025]{ctim-rover}
Lindenbauer, M.; et al.
\newblock From knowledge to noise: {CTIM}-Rover and the pitfalls of episodic memory in software engineering agents.
\newblock In \emph{Proceedings of REALM (ACL Workshop)}, 2025.

\bibitem[Lingam et al., 2024]{dot-bank}
Lingam, V.; et al.
\newblock {DoT-Bank}: Diverse self-reflections and exemplar trajectories for agent memory.
\newblock \emph{arXiv preprint}, 2024.

\bibitem[Agent Workflow Memory]{awm}
Agent Workflow Memory.
\newblock \emph{arXiv preprint}, 2024.

\bibitem[Packer et al., 2023]{memgpt}
Packer, C.; Fang, V.; Patil, S. G.; Lin, K.; Wooders, S.; and Gonzalez, J. E.
\newblock {MemGPT}: Towards {LLMs} as operating systems.
\newblock \emph{arXiv preprint arXiv:2310.08560}, 2023.

\bibitem[Park et al., 2023]{generative-agents}
Park, J. S.; O'Brien, J. C.; Cai, C. J.; Morris, M. R.; Liang, P.; and Bernstein, M. S.
\newblock Generative agents: Interactive simulacra of human behavior.
\newblock In \emph{Proceedings of UIST}, 2023.

\bibitem[Zhang et al., 2025]{g-memory}
Zhang, Y.; et al.
\newblock {G-Memory}: Tracing hierarchical memory for multi-agent systems.
\newblock In \emph{NeurIPS}, 2025.

\bibitem[Shah et al., 2025]{evolve-mem}
Shah, R.; et al.
\newblock {EVOLVE-MEM}: A self-adaptive hierarchical memory architecture.
\newblock In \emph{NeurIPS}, 2025.

\bibitem[H-MEM]{h-mem}
{H-MEM}: Hierarchical memory for high-efficiency long-term reasoning.
\newblock \emph{arXiv preprint arXiv:2507.22925}, 2025.

\bibitem[Context-Picker]{context-picker}
{Context-Picker}: Reasoning-aware variable-length evidence selection.
\newblock \emph{arXiv preprint arXiv:2512.14465}, 2025.

\bibitem[MEM1]{mem1}
{MEM1}: Learning to synergize memory and reasoning.
\newblock In \emph{NeurIPS}, 2025.

\bibitem[MemFly]{memfly}
{MemFly}: On-the-fly memory optimization via information bottleneck.
\newblock \emph{arXiv preprint arXiv:2602.07885}, 2025.

\bibitem[DCPM]{dcpm}
Memory beyond recall: A dual-process cognitive memory system.
\newblock \emph{arXiv preprint arXiv:2606.09483}, 2025.

\bibitem[CDMem]{cdmem}
Context-dependent memory for {LLM} agents.
\newblock In \emph{Proceedings of NAACL}, 2025.

\bibitem[Tulving, 1972]{tulving-episodic}
Tulving, E.
\newblock Episodic and semantic memory.
\newblock In \emph{Organization of Memory}, Academic Press, 1972.

\bibitem[Tulving and Thomson, 1973]{tulving-encoding}
Tulving, E.; and Thomson, D. M.
\newblock Encoding specificity and retrieval processes in episodic memory.
\newblock \emph{Psychological Review}, 80(5):352--373, 1973.

\bibitem[Schank, 1982]{schank-dynamic}
Schank, R. C.
\newblock \emph{Dynamic Memory: A Theory of Reminding and Learning in Computers and People}.
\newblock Cambridge University Press, 1982.

\bibitem[Sweller, 1988]{cognitive-load}
Sweller, J.
\newblock Cognitive load during problem solving: Effects on learning.
\newblock \emph{Cognitive Science}, 12(2):257--285, 1988.

\bibitem[Reimers and Gurevych, 2019]{sentence-transformers}
Reimers, N.; and Gurevych, I.
\newblock {Sentence-BERT}: Sentence embeddings using {Siamese} {BERT}-networks.
\newblock In \emph{Proceedings of EMNLP-IJCNLP}, 2019.

\bibitem[Jimenez et al., 2024]{swe-bench}
Jimenez, C. E.; Yang, J.; Wettig, A.; Yao, S.; Pei, K.; Press, O.; and Narasimhan, K.
\newblock {SWE}-bench: Can language models resolve real-world {GitHub} issues?
\newblock In \emph{ICLR}, 2024.

\bibitem[Feng et al., 2020]{codebert}
Feng, Z.; Guo, D.; Tang, D.; Duan, N.; Feng, X.; Gong, M.; Shou, L.; Qin, B.; Liu, T.; Jiang, D.; and Zhou, M.
\newblock {CodeBERT}: A pre-trained model for programming and natural languages.
\newblock In \emph{Findings of EMNLP}, 2020.

\bibitem[Wang et al., 2021]{codet5}
Wang, Y.; Wang, W.; Joty, S.; and Hoi, S. C.
\newblock {CodeT5}: Identifier-aware unified pre-trained encoder-decoder models.
\newblock In \emph{Proceedings of EMNLP}, 2021.

\bibitem[Roziere et al., 2023]{code-llama}
Rozi\`{e}re, B.; et al.
\newblock {Code Llama}: Open foundation models for code.
\newblock \emph{arXiv preprint arXiv:2308.12950}, 2023.

\bibitem[Brown et al., 2020]{icl}
Brown, T.; et al.
\newblock Language models are few-shot learners.
\newblock In \emph{NeurIPS}, 2020.

\bibitem[ActMem, 2025]{actmem}
{ActMem}: Bridging the gap between memory retrieval and reasoning in {LLM} agents.
\newblock \emph{arXiv preprint arXiv:2603.00026}, 2025.

\end{thebibliography}

\end{document}
